\section{Receiver-Calibrated Prototype Adaptation}
\label{sec:rcpa}

\subsection{Frozen backbone}
We use the OOB-guided RF-HSTU hybrid model from our related submitted work as a fixed feature extractor~\cite{paper1_placeholder}.
No backbone parameters are updated in this paper.
All calibration operates on fused embeddings $z(w)$ extracted from target-receiver windows.
This design isolates the value of target-receiver calibration from any architecture change.

\subsection{Baselines}
\textbf{Source classifier.}
The source-trained softmax head $g(\cdot)$ is applied directly to target-query embeddings:
$\hat{y}=\arg\max_c g(z)_c$.

\textbf{Source prototype (RCPA-S).}
Source-receiver class prototypes are
\begin{equation}
\mu_s^{(c)}=\frac{1}{|\mathcal{S}_s^{(c)}|}\sum_{w\in\mathcal{S}_s^{(c)}} z(w),
\end{equation}
where $\mathcal{S}_s^{(c)}$ is the set of labeled source windows for class $c$.
Target queries are classified by minimum cosine distance to $\{\mu_s^{(c)}\}$, following the prototypical classification principle~\cite{snell2017prototypical}.

\subsection{Target-receiver prototype calibration}
For each device class $c$, we form a target prototype from labeled support windows on $r_t$:
\begin{equation}
\mu_t^{(c)}=\frac{1}{K}\sum_{w\in\mathcal{S}_t^{(c)}} z(w),
\end{equation}
where $\mathcal{S}_t^{(c)}$ contains $K$ labeled windows sampled uniformly from support Block~B for device $c$.
Each query device file is represented by the mean fused embedding of its query-block windows $\bar{z}=\mathrm{mean}_{w\in\mathcal{Q}} z(w)$ and classified by
\begin{equation}
\hat{y}=\arg\min_c \; d_{\cos}(\bar{z},\mu^{(c)}).
\end{equation}

\textbf{RCPA-T (primary).}
Use target prototypes only: $\mu^{(c)}=\mu_t^{(c)}$.

\textbf{RCPA-B (ablation).}
Fixed source-target blend
\begin{equation}
\mu^{(c)}=\alpha\mu_s^{(c)}+(1-\alpha)\mu_t^{(c)}, \quad \alpha=0.5,
\end{equation}
included as ablation only; not tuned on query data.

\subsection{Block-disjoint split protocol}
Each target file contributes 256 deterministic windows partitioned into four contiguous blocks of 64 windows:
Block~A (calibration, unlabeled), Block~B (support, labeled), Blocks~C+D (query, evaluation only).
Three split repeats rotate block roles so that each block serves each role once across repeats.
We assert zero overlap among calibration, support, and query window sets per device.

Because each device has one capture file per receiver, $K$ always denotes labeled \emph{windows} per device rather than independent files.
For example, $K=5$ uses five labeled windows from Block~B, not five separate enrollment captures.

\subsection{Leakage prevention}
Query windows are never used for prototype construction, pseudo-label selection, confidence thresholding, or equalization-statistic estimation.
Support labels are used only for RCPA-T/RCPA-B prototype construction.
Unsupervised adaptation baselines may read Block~A only; evaluation uses Blocks~C+D only.

\subsection{Unsupervised target calibration baselines}
\textbf{Entropy-min TTA.}
Following TENT~\cite{wang2021tent}, we update batch-norm or lightweight adapter statistics on unlabeled Block~A windows by minimizing prediction entropy, then evaluate on query blocks.

\textbf{Pseudo-proto TTA.}
High-confidence pseudo-labeled Block~A windows update class prototypes online, following source-free and pseudo-label adaptation practice~\cite{liang2020shot,yang2025scrffi}; we sweep confidence thresholds $\{0.5,0.7,0.8,0.9,0.95\}$ in appendix experiments.

\subsection{OOB representation equalization (auxiliary)}
As mechanism validation, we apply embedding-level CORAL alignment~\cite{sun2016deepcoral} to reduce OOB receiver probe accuracy before RCPA evaluation.
Statistics are estimated from source embeddings and target calibration Block~A only, following unsupervised domain-adaptation practice~\cite{ben2010domain_adaptation,sun2016deepcoral}.
This analysis tests whether suppressing receiver information in OOB embeddings alone restores classification without labeled target prototypes.
It is not proposed as a primary deployment method.

\subsection{K-shot parametric baselines (same protocol)}
For comparison, we also evaluate K-shot multinomial logistic regression and classifier-head fine-tuning on frozen embeddings under the same support/query split.
These baselines represent strong parametric target calibration alternatives to non-parametric prototype voting~\cite{snell2017prototypical,wen2016center_loss}.

\subsection{Deployment modes}
\begin{itemize}
  \item \textbf{Source-only:} no target labels; source classifier or RCPA-S.
  \item \textbf{Unlabeled target calibration:} entropy-min or pseudo-proto TTA on Block~A only.
  \item \textbf{$K$-shot labeled target calibration:} RCPA-T using Block~B support windows.
\end{itemize}

\textbf{Important:} RCPA-T is post-hoc target-receiver calibration, not a new RF-HSTU architecture and not a backbone retraining method.

\begin{figure}[t]
  \centering
  \includegraphics[width=\linewidth]{fig2_rcpa_shotcurve}
  \caption{RCPA-T file-level accuracy versus $K$ labeled calibration windows per device (frozen full-mode results).}
  \label{fig:rcpa_shotcurve}
\end{figure}
